% !TEX root = ../main.tex

\chapter{Figures and Tables}
\label{ch:examples}

This chapter is intentionally practical: it contains the small and full-width
float patterns that most thesis projects need. The figures use copied example
images from the source thesis so the template shows the actual supported image
workflow.

\section{Figures}
\label{sec:figures}

\begin{figure}
	\centering
	\includegraphics[width=0.86\linewidth]{tabular/social_choice_borda_bt.pdf}
	\caption{A regular figure. Its caption sits in the margin in the standard
	kaobook layout. Use this for figures that fit naturally in the text column.}
	\label{fig:regular-figure}
\end{figure}

Reference regular figures with \texttt{\string\ref}: Figure~\ref{fig:regular-figure}
is a text-column float. With an image file, use \texttt{\string\includegraphics}
with a width such as \texttt{\string\textwidth} or \texttt{0.8\string\linewidth}.

\begin{figure*}
	\centering
	\includegraphics[width=0.92\linewidth]{tabular/social_choice_borda_bt.pdf}
	\caption{A full-width figure. Use \texttt{figure*} when the content needs
	the text column and the margin together.}
	\label{fig:full-width-figure}
\end{figure*}

Figure~\ref{fig:full-width-figure} uses the full available page width. This is
useful for multi-panel figures, dense plots, or diagrams with long labels.

\section{Tables}
\label{sec:tables}

\begin{table}
	\centering
	\begin{tabular}{lrr}
		\toprule
		Method & Score & Time \\
		\midrule
		Baseline & 0.71 & 12.4 \\
		Proposed & 0.83 & 10.9 \\
		\bottomrule
	\end{tabular}
	\caption{A regular table for compact results. Keep column names short and
	put details in the caption or surrounding text.}
	\label{tab:regular-table}
\end{table}

Table~\ref{tab:regular-table} is a normal text-column table.

\begin{table*}
	\centering
	\small
	\begin{tabularx}{0.95\linewidth}{lrrrrX}
		\toprule
		Setting & Train & Valid & Test & Runtime & Notes \\
		\midrule
		Small & 100 & 20 & 20 & 1.2 min & Useful for smoke tests and examples. \\
		Medium & 1000 & 200 & 200 & 8.4 min & Fits comfortably in a wider table. \\
		Large & 10000 & 2000 & 2000 & 51.0 min & Use wide tables for many columns or longer notes. \\
		\bottomrule
	\end{tabularx}
	\caption{A full-width table. Use \texttt{table*} for wide result summaries
	or tables with explanatory columns.}
	\label{tab:full-width-table}
\end{table*}

Table~\ref{tab:full-width-table} demonstrates the wide-table layout and
\texttt{tabularx} for a flexible notes column.
